\section{Construction primitives}

The algebraic language should include direct sums, tensor products,
transposes and involutions, block composition, quotient modules, symmetry
actions, and explicit choices of basis. Not every operation applies to
every dialect. A construction rule should record its hypotheses and a
verifier for its finite output.

\subsection{Hypergraph product}
For binary matrices $A$ and $B$ of shapes $m_1\times n_1$ and
$m_2\times n_2$, define
\begin{align}
 H_X&=[A\otimes I_{n_2}\mid I_{m_1}\otimes B^T],\\
 H_Z&=[I_{n_1}\otimes B\mid A^T\otimes I_{m_2}].
\end{align}
Then $H_XH_Z^T=A\otimes B^T+A\otimes B^T=0$ over $\F$.
The blocklength is $n_1n_2+m_1m_2$. Writing
$k_i=n_i-\rank H_i$ and $k_i^T=m_i-\rank H_i$, the logical count is
$k_1k_2+k_1^Tk_2^T$~\cite{tillich2014}. The ability to state such a
formula at construction level is a benefit of retaining structure rather
than only its expanded parity-check matrix.

\subsection{Commuting blocks and lifts}
If square binary matrices $A$ and $B$ commute, then
\begin{equation}
 H_X=[A\mid B],\qquad H_Z=[B^T\mid A^T]
\end{equation}
form a CSS code because $AB+BA=0$. Commuting shift operators, represented
by polynomials in a finite quotient ring, provide compact parameters for
such blocks. Bivariate bicycle codes illustrate the practical interest
of this design space~\cite{bravyi2024}.

A symbolic primitive should define its coefficient ring, quotient
relations, involution, basis ordering, and expansion map. For
noncommutative group algebras, left and right module actions cannot be
silently interchanged. Lifted-product constructions exploit this richer
structure, and asymptotically good qLDPC families are already known
\cite{panteleev2022}. Merely validating commutation does not reproduce
their distance proofs.

\subsection{A shared abelian lift}
For $A\in R^{m_A\times n_A}$ and $B\in R^{m_B\times n_B}$,
$R=\F[G]$ with $G$ finite abelian, let star combine matrix transpose
and inversion of group monomials. Define
\begin{align}
 H_X&=[A\otimes I_{n_B}\mid I_{m_A}\otimes B^*],\\
 H_Z&=[I_{n_A}\otimes B\mid A^*\otimes I_{m_B}].
\end{align}
This uses the lifted-product convention of~\cite[Section III-D]{panteleev2021}
with its second seed replaced by $B^*$. Commutativity gives
$H_XH_Z^*=A\otimes B^*+A\otimes B^*=0$. Regular representation
preserves products and sends star to binary transpose, proving CSS
validity after expansion.

For $L=|G|$, the length is $L(n_An_B+m_Am_B)$; independently expanding
both seeds before taking an ordinary hypergraph product instead costs
$L^2(n_An_B+m_Am_B)$ qubits. No generic field-rank dimension formula is
assumed over this ring. Uniform bounds on seed row and column sums of
monomial support sizes control expanded locality. General noncommutative
extensions require compatible commuting actions and are outside this
construction's scope.

\subsection{Families and locality}
A qLDPC claim concerns a family with generator weight and participation
per physical qubit bounded independently of blocklength. A sparse-looking
single matrix is insufficient. Retain the original presentation, the
family parameter, and a uniform bound. Geometric locality imposes an
additional embedding constraint; bounded check weight alone does not
make an interaction geometrically local.

Homological algebra is the general framework here. Algebraic geometry,
expander graphs, group theory, and topology supply distinct construction
resources. No assumption that all useful codes arise from geometric
surfaces is required~\cite{breuckmann2021}.
